\documentclass[10pt,a4paper]{article}
\usepackage[utf8]{inputenc}
\usepackage[a4paper,left=3cm,right=2cm]{geometry}
\usepackage{amsmath}
\usepackage[thmmarks, amsmath, thref]{ntheorem}
\usepackage{amsfonts}
\usepackage{amssymb}
\usepackage{fancyhdr}
\pagestyle{fancy}
\usepackage{anysize}
\usepackage{graphicx}
\usepackage{float}
\usepackage{hyperref}
\usepackage{enumitem}
\newtheorem{theorem}{Theorem}
\theoremsymbol{\ensuremath{\blacksquare}}
\newtheorem*{solution}{Solución:}
\usepackage{float}
\usepackage{xcolor}

% Margenes y formalidades

\textwidth 6.25in % Margen Derecho y ancho del texto
\textheight 8.5in
\oddsidemargin 0in % Margen Izquierdo
\headheight 0in % Largo del texto
\leftmargin 1in
\parindent 0em % Salto entre la indentación (Sangría)https://www.overleaf.com/4519127669sdpptkhvsbfn
\topmargin -0.5in % Margen Superior
\parskip 0.5ex % Salto entre parrafos
\baselineskip 1pt




%\chead{}                                      % texto centro de la cabecera
\rhead{Universidad Técnica Federico Santa María\\Ingeniería Comercial } % número de página a la derecha
\renewcommand{\headrulewidth}{0.4pt} % grosor de la línea de la cabecera
\renewcommand{\footrulewidth}{0.4pt} % grosor de la línea del pie

\begin{document}
\vspace{5cm}



\begin{center}
\begin{large}
\textbf{\\CERTAMEN 3}
\end{large}\\
\end{center}

\begin{center}
\begin{large}
\textbf{Econometría ICS-294}
\end{large}\\
\vspace{0.2cm}

\end{center}

\textbf{Instrucciones:} 

\begin{itemize}

    \item Usted dispone de 60 minutos  para intentar los 100 puntos del certamen.
    \item Cuide la presentación y redacción de sus preguntas.
    \item No puede utilizar celular, computador o tablet, pero sí una calculadora científica.
     
\end{itemize}

\section{Preguntas Cortas (30 Puntos)}

\begin{enumerate}

\item (15 Puntos) ¿En qué casos, y por qué motivos, conviene transformar las variables a primeras diferencias cuando se trabaja con series de tiempo?

\item (15 Puntos) Defina claramente las dos condiciones que debe cumplir una variable instrumental para ser considerada válida e indique qué problema aparece si se viola cada una de ellas.  

\end{enumerate}

\section{Ejercicios (70 Puntos)}

\begin{enumerate}

\item (35 Puntos) Se dispone de una muestra de 350 observaciones con la cual se estimó el siguiente modelo de regresión lineal múltiple:
$$Salario_i = \beta_0 + \beta_1\,\textit{Educ}_i + \beta_2\,\textit{Exper}_i + u_i$$
Se desea contrastar el supuesto de heterocedasticidad y, para ello, se aplicará el método de White simplificado.
\begin{enumerate}
    \item (10 Puntos) Plantee la regresión auxiliar para el método de White simplificado.
    
    
    \item (10 puntos) Formule la hipótesis nula y la hipótesis alternativa para contrastar la presencia de heterocedasticidad.
    
    \item (15 Puntos) Si el $R^2$ de la regresión auxiliar es $0.14$, contraste, al 5\% de significación, la existencia de heterocedasticidad mediante el estadístico del multiplicador de Lagrange.
    
\end{enumerate}



\item (35 Puntos) Se sospecha que la función de producción del periodo t ($y_t$) sigue una serie de tiempo de tendencia estocástica; por lo tanto, se plantea la siguiente regresión:
$$\Delta y_t = \beta_0 + \phi y_{t-1} + v_t$$
Y se encuentran los siguientes resultados para una muestra de 200 datos:
\begin{center}
    \begin{figure}[H]
        \centering
        \includegraphics[width=0.6\linewidth]{certamenes/Resultado Serie de tiempo.jpg}
        \caption{Resultados Regresión}
        \label{fig:enter-label}
    \end{figure}
\end{center}
\begin{enumerate}
    \item (10 Puntos) Plantee la hipotesis nula y la hipotesis alternativa que permite testear si la regresión en diferencias tiene raiz unitaria.
    
    \item (20 Puntos) Testee con un 5\% de significancia si la regresión en primeras diferencias tiene raiz unitaria y concluya. 
    
\end{enumerate}




\end{enumerate}


\section*{Formulario}
\begin{itemize}

\item Tabla $\chi^2$
\begin{center}
    \begin{tabular}{|c|c|c|c|}
    \hline
        df & 0.90 & 0.95 & 0.90 \\ 
        \hline
        1 &  6.635 & 3.841 & 2.706\\
        \hline
        2 &  9.210 & 5.991 & 4.605 \\
        \hline
        3 & 11.345 & 7.815 & 6.251  \\
        \hline
        4 &  13.277 & 9.488 & 7.779 \\
        \hline
        5 &  15.086 & 11.070 & 9.236 \\
        \hline
        6 &  16.812 & 12.592 & 10645 \\
        \hline
    \end{tabular}
\end{center}

\item Multiplicado de Lagrange

$$\color{red}{LM = nR^2 \sim \chi^2}$$

\item Test F para regresión auxiliar.
\begin{center}
 {  \color{red}     \(F=\dfrac{R_{\text{aux}}^{2}/q}{(1-R_{\text{aux}}^{2})/(n-k-1)}
          \sim F_{q,\;n-k-1}\).}
\end{center}



\item Tabla de test D-F
\begin{center}
    \begin{figure}[H]
        \centering
        \includegraphics[width=0.7\linewidth]{sem14/tabla.png}
    \end{figure}
\end{center}

\end{itemize}


\end{document}